\section{Introduction}\label{intro:section}

For three points $p,q,r\in\mathbb R^2$, write
\[
 \operatorname{Area}(pqr)=\frac12|\det(q-p,r-p)|.
\]
If $P\subset[0,1]^2$ is finite and $|P|\ge3$, define
\[
 \Delta(P)=\min_{\{p,q,r\}\in\binom P3}\operatorname{Area}(pqr),
 \qquad
 \Delta(n)=\max_{\substack{P\subset[0,1]^2\\|P|=n}}\Delta(P).
\]
The minimum is over all unordered triples of distinct points. In particular,
a collinear triple makes $\Delta(P)=0$.
The maximum defining $\Delta(n)$ exists: the minimum area is continuous on
the compact space of ordered $n$-tuples, it is positive for suitable points
on a parabola, and a tuple with a repeated point has minimum area zero.

The original Heilbronn conjecture asked whether $\Delta(n)=O(n^{-2})$;
see Roth~\cite[p.~198]{Roth1951}. Two elementary constructions attain this
scale. Erd\H{o}s's finite-field parabola, recorded in
Roth~\cite[Appendix]{Roth1951}, gives a grid configuration whose triangle
determinants are nonzero integers before scaling. Random sampling followed
by deleting one point from each small triangle also gives
$\Delta(n)\gg n^{-2}$. Koml\'os, Pintz and Szemer\'edi improved this to
$\Delta(n)\gg(\log n)/n^2$, disproving the original
conjecture~\cite{KPS1982}. Their argument represents the small-area triples
in a random point set as edges of a hypergraph, removes short cycles,
and uses a stronger independence bound to select a large subset containing
none of those triples~\cite[Section~2]{KPS1982}.

The almost $n^{-2}$ formulation, discussed in
Zakharov's survey~\cite[Section~3]{Zakharov2026}, asks whether, for every
$\varepsilon>0$, there are constants $C_\varepsilon$ and
$n_0(\varepsilon)$ such that
\begin{equation}\label{intro:almost}
 \Delta(n)\le C_\varepsilon n^{-2+\varepsilon}
 \qquad(n\ge n_0(\varepsilon)).
\end{equation}
We disprove this formulation by a power improvement in the lower bound.

\begin{theorem}\label{intro:main}
There are absolute constants $\eta,c_1>0$ and an integer $n_0\ge3$ such that
\[
 \Delta(n)\ge c_1n^{-2+\eta}
 \qquad\text{for every integer }n\ge n_0.
\]
\end{theorem}

In particular, $\Delta(n)\ge n^{-2+\eta/2}$ for every sufficiently large
$n$. The ratio of the theorem's lower bound to $n^{-2+\eta/2}$ is
$c_1n^{\eta/2}\to\infty$, which contradicts \eqref{intro:almost} at
$\varepsilon=\eta/2$. Section~\ref{alt:section} gives an explicit, though
extremely small, admissible exponent.

On the upper-bound side, Roth~\cite{Roth1951} proved the first
$o(n^{-1})$ estimate by a density-increment argument, followed by refinements
of Schmidt~\cite{Schmidt1972} and Roth~\cite{Roth1972a,Roth1972b}.
Koml\'os, Pintz and Szemer\'edi~\cite{KPS1981} refined Roth's analytic method
to obtain $\Delta(n)\ll_\varepsilon n^{-8/7+\varepsilon}$.
Cohen, Pohoata and Zakharov developed new incidence-geometric approaches,
first obtaining $\Delta(n)\le n^{-8/7-1/2000}$ for sufficiently large
$n$~\cite[Theorem~1.1]{CPZ2023}, and subsequently proving
$\Delta(n)\ll_\varepsilon n^{-7/6+\varepsilon}$ for every
$\varepsilon>0$~\cite[Theorem~1.8]{CPZ2025}.

Earlier preprints claim stronger power lower bounds than ours, with
unresolved issues in the following versions. Ellmann's version~12
acknowledges an unproved local uniformity assumption and describes the
argument as heuristic~\cite[p.~3]{Ellmann2025}; the deletion estimate uses
this assumption~\cite[Theorem~4, pp.~5--6]{Ellmann2025}.
In Agama's version~13, Theorem~4.1~\cite[pp.~12--13]{Agama2026},
counting the center with the boundary points introduces a collinear triple
consisting of the two endpoints of a diameter and their midpoint.
Omitting the center leaves the every-triple estimate unproved, since the
argument estimates triangles formed by the center and adjacent boundary
points. These objections concern the cited arguments, not the impossibility
of their asserted bounds.

\subsection{The two congruence conditions}

The proof works first with integer columns $u=(u_1,u_2,u_3)^{\mathsf T}$
in a box of the form
\[
 0\le u_1,u_2<N,\qquad N\le u_3<2N.
\]
Such a column represents the unit-square point
$\pi(u)=(u_1/u_3,u_2/u_3)$. Three columns forming a matrix $A$ give a
triangle of area
\begin{equation}\label{intro:area}
 \frac{|\det A|}{2A_{31}A_{32}A_{33}}.
\end{equation}
We therefore seek many columns for which proportional pairs and triples
with small absolute determinant are sufficiently rare to delete.

The first congruence condition separates most triples before they are
lifted to integers. Give each column an independent uniform label
$\xi\in\mathbb F_{r^d}$, where $d$ is fixed and odd and the prime $r$ grows.
The columns $(1,\xi,\xi^2)^{\mathsf T}$
have nonzero determinant whenever their labels are distinct. This is the
finite-field parabola underlying Erd\H{o}s's construction recorded in
Roth~\cite[Appendix]{Roth1951}. Taking the
field norm to $\mathbb F_r$ produces an alternating polynomial in three groups of
coordinates. We express that polynomial as a sum of monomial determinants
and encode it in one designated coefficient of a determinant whose entries
are base-$B$ expansions.
Large-base control of carries has classical precedents in the
progression-free constructions of Salem and Spencer~\cite{SalemSpencer1942}
and Behrend~\cite{Behrend1946}; here it isolates a determinant coefficient.
Modulo $h=B^k$, with $k$ fixed, distinct labels then preclude a small
determinant. Independent choices of digits with the prescribed residues
leave enough randomness even when labels coincide.

A common random matrix of determinant one mixes the three rows modulo
$h$ without changing their determinant. Conditional on the digit columns,
the resulting matrix is uniform on a special-linear orbit. The rows of
every integral lift lie in a lattice whose index, together with the orbit
size, can be read from a diagonal form.
An anisotropic lattice count controls lifts of any fixed nonzero
determinant, while a conditional moment bound handles singular digit
matrices and repeated labels.

The fixed nonzero-determinant count does not cover determinant zero, where
all three rows can lie in one integral plane. We control this case with a
second congruence condition. At an independent prime $q$, choose residues
from a small portion of an affine quadratic surface with no three points on a line. A carefully restricted random
translation excludes short integer relations; a random invertible linear
map permits uniform counting. The Chinese remainder theorem combines the
two congruence conditions, and uniform lifting supplies integer columns
in the box.

For a singular lifted matrix with pairwise distinct projected columns,
a primitive integer null vector determines a rank-two row lattice. We sum over those
vectors, treating equal residue columns and low rank modulo $q$
separately. This supplies the missing count for collinear projected triples.

The exceptional triples require repeated labels, whose probability
provides the saving used in an elementary deletion argument. Thus the
improvement comes from the arithmetic distribution of sampled points.
Formula~\eqref{intro:area} then transfers the determinant bound to every
triangle of the retained point set.

\subsection{Organization and conventions}

Section~\ref{lat:section} proves the lattice estimates.
Section~\ref{norm:section} constructs the norm polynomial and digit
distribution, and Section~\ref{orb:section} estimates its matrix orbits.
Section~\ref{aux:section} constructs the auxiliary residues.
Sections~\ref{sam:section} and~\ref{zero:section} count small determinants,
including zero. Section~\ref{alt:section} completes the deletion and scaling
argument and passes to every sufficiently large cardinality.
Appendix~\ref{app:prime-interval} gives the elementary prime-interval proof
used in the parameter choices and interpolation.

Lengths are Euclidean. A primitive integer vector has coordinates with
greatest common divisor one. The covolume of a lattice is measured in its
real span. The notation $F\ll G$ means $|F|\le CG$; constants may depend on
the fixed construction parameters, but never on the growing prime or on
the sampled labels and columns. All logarithms are natural.
